\documentclass[11pt,a4paper]{article}

\usepackage[utf8]{inputenc}
\usepackage[T1]{fontenc}
\usepackage{geometry}
\usepackage{hyperref}
\usepackage{enumitem}
\usepackage{xcolor}
\usepackage{fancyhdr}
\usepackage{amsmath}
\usepackage{booktabs}
\usepackage{listings}

\geometry{margin=1in}
\hypersetup{
  colorlinks=true,
  linkcolor=blue,
  urlcolor=blue
}

\setlist{itemsep=0.35em, topsep=0.35em}

\pagestyle{fancy}
\fancyhf{}
\lhead{Object Oriented Programming (B.Tech CSE)}
\rhead{Unit 02: Classes, Inheritance, Packages and Interfaces}
\cfoot{\thepage}
\setlength{\headheight}{14pt}

\lstset{
  language=Java,
  basicstyle=\ttfamily\small,
  keywordstyle=\color{blue},
  commentstyle=\color{gray},
  stringstyle=\color{teal},
  showstringspaces=false,
  columns=fullflexible,
  frame=single
}

\begin{document}

\begin{center}
  {\LARGE \textbf{Assignment -- Unit 02}}\\[0.25em]
  {\large Object Oriented Programming (CSEG1044)}\\[0.25em]
  \normalsize (Classes/objects; constructors; \texttt{this}; overloading; \texttt{static}; access modifiers; packages; interfaces; inheritance; abstract; \texttt{final})
\end{center}

\noindent
\textbf{Instructor:} Tofik Ali \hfill \textbf{Submission Deadline:} March 31, 2026\\
\textbf{Student Name:} \_\_\_\_\_\_\_\_\_\_\_\_\_\_\_\_\_ \hfill \textbf{Roll No.:} \_\_\_\_\_\_\_\_\\

\vspace{0.6em}
\hrule
\vspace{0.8em}

\textbf{Instructions}
\begin{itemize}[leftmargin=*]
  \item Answer \textbf{all} questions. Write assumptions where needed.
  \item There is \textbf{no time limit}. Submit on or before \textbf{March 31, 2026}.
  \item For programming questions: write clean Java code, enforce invariants, and show a sample run/output.
  \item Unless specified, use only standard Java libraries.
  \item If a question requires packages, follow correct folder structure (package name $\leftrightarrow$ directory path).
\end{itemize}

\vspace{0.6em}

\section*{Easy (10 Questions)}
\begin{enumerate}[label=E\arabic*., leftmargin=*]
  \item \textbf{Class vs object (short):} Write 3 differences between a class and an object, with one small example.

  \item \textbf{Encapsulation (fix design):} Given a class with public fields, rewrite it so fields become \texttt{private}
    and updates happen via methods that validate inputs. Explain in 4--6 lines why this is safer.

  \item \textbf{Constructors:} What is a constructor? Explain default constructor vs parameterized constructor (4--6 lines).

  \item \textbf{\texttt{this} vs \texttt{this(...)}:} Explain the difference and write one code snippet for each.

  \item \textbf{Overloading rules:} Give 2 valid overloaded method signatures and 1 invalid example
    (invalid = changing only return type).

  \item \textbf{\texttt{static} basics:} Explain \texttt{static} fields and methods with one real example (e.g., counter/utility).

  \item \textbf{Access modifiers (quick):} For each modifier, write where it is visible:
    \texttt{private}, default (package-private), \texttt{protected}, \texttt{public}.

  \item \textbf{Packages (naming):} Write a good package name for a project called ``Campus Resource Manager''
    and show the matching folder path.

  \item \textbf{Interface (idea):} Define what an interface is in 3--5 lines. Write a tiny interface with one method
    and one class that implements it.

  \item \textbf{is-a vs has-a:} For each pair, decide is-a or has-a and justify in 1--2 lines:
    \textbf{(a)} Car--Engine \textbf{(b)} Manager--Employee \textbf{(c)} Student--IDCard.
\end{enumerate}

\section*{Medium (10 Questions)}
\begin{enumerate}[label=M\arabic*., leftmargin=*]
  \item \textbf{Student class (invariant):} Implement \texttt{Student} with private fields:
    \texttt{rollNo}, \texttt{name}, \texttt{cgpa}. Enforce invariant \texttt{0 <= cgpa <= 10} using validation.
    Provide methods to update CGPA safely (no direct setter without validation).

  \item \textbf{Employee with constructor chaining + static counter:} Implement \texttt{Employee} with fields:
    \texttt{id}, \texttt{name}, \texttt{salary}. Generate \texttt{id} automatically using a static counter.
    Include at least two constructors using \texttt{this(...)} chaining.

  \item \textbf{BankAccount mini-design:} Implement \texttt{BankAccount} with:
    account number, balance, \texttt{deposit}, \texttt{withdraw}, \texttt{transferTo}.
    Add at least one overloaded constructor (e.g., opening balance optional).

  \item \textbf{Packages (multi-file):} Create these files and packages:
    \begin{itemize}[leftmargin=*]
      \item \texttt{com.upes.crm.model.Resource}
      \item \texttt{com.upes.crm.app.Main}
    \end{itemize}
    In \texttt{Main}, create a \texttt{Resource} object and print it.
    (Override \texttt{toString()} in \texttt{Resource}.)

  \item \textbf{Interface + polymorphism:} Define \texttt{Payable} with method \texttt{double amountDue()}.
    Implement it in \texttt{Invoice} and \texttt{SalarySlip}. In \texttt{main}, store them in a \texttt{Payable[]} array and compute total.

  \item \textbf{Inheritance + overriding:} Create \texttt{Shape} with method \texttt{double area()}.
    Create \texttt{Circle} and \texttt{Rectangle} subclasses that override \texttt{area()}.
    Demonstrate dynamic dispatch using a \texttt{Shape} reference.

  \item \textbf{Abstract class:} Create an abstract class \texttt{EmployeeBase} with abstract method
    \texttt{double calculateSalary()}. Create two subclasses (e.g., \texttt{FullTime}, \texttt{Intern}) and implement salary logic.

  \item \textbf{\texttt{final} keyword:} Demonstrate:
    \begin{enumerate}[label=(\alph*), leftmargin=*]
      \item a \texttt{final} variable,
      \item a \texttt{final} method,
      \item a \texttt{final} class.
    \end{enumerate}
    Explain the design intention in 6--10 lines.

  \item \textbf{Access modifiers across packages:} Create a base class in \texttt{pkg1} with:
    one \texttt{protected} method, one package-private method (no keyword), and one \texttt{public} method.
    Create a subclass in \texttt{pkg2} and show which methods can/can't be accessed (compile-time).

  \item \textbf{Built-in interface: Comparable:} Implement \texttt{Comparable<Student>} in \texttt{Student}
    (choose a natural ordering: roll number or cgpa).
    Create a \texttt{Student[]} array and sort it using \texttt{Arrays.sort}.
\end{enumerate}

\section*{Hard (10 Questions)}
\begin{enumerate}[label=H\arabic*., leftmargin=*]
  \item \textbf{Mini system (design + code):} Implement a very small ``Campus Resource Manager'' core using classes:
    \texttt{Resource(id,name,type,capacity)} and \texttt{Booking(resourceId,userName,start,end)}.
    Add basic validations (capacity $>$ 0, start $<$ end). (No DB/UI needed.)

  \item \textbf{Interface-driven design:} Create an interface \texttt{Storage} with methods:
    \texttt{save(String key, String value)} and \texttt{String find(String key)}.
    Implement two classes:
    \texttt{InMemoryStorage} (use arrays) and \texttt{FileStorage} (you may keep file part as TODO comment).
    Write a \texttt{main} that works with \texttt{Storage} reference.

  \item \textbf{Composition example:} Implement \texttt{Car} has-a \texttt{Engine}.
    Engine should have fields like \texttt{hp} and \texttt{type}. Show that \texttt{Car} delegates to \texttt{Engine}.
    Write 5--8 lines: why composition is preferred here.

  \item \textbf{Inheritance with \texttt{super}:} Create \texttt{Person(name)} and \texttt{Student(name,rollNo)}.
    Use \texttt{super(name)} correctly. Override \texttt{toString()} and print both.

  \item \textbf{Overriding vs overloading:} Create a short example that shows both in the same codebase.
    Explain in 6--10 lines how Java chooses which method to call.

  \item \textbf{Abstract class vs interface (design decision):} Pick one scenario (payment, logging, notification, or sorting)
    and argue (10--14 lines) whether you would use an interface or an abstract class. Include one code sketch.

  \item \textbf{Access control (API design):} Design a package \texttt{com.upes.bank} with:
    a public class \texttt{BankAccount} and a package-private helper class \texttt{Validator}.
    Explain why hiding helper classes helps maintainability.

  \item \textbf{Small library (search overloading):} Create classes \texttt{Book(title,author,year)} and \texttt{Library}.
    In \texttt{Library}, implement overloaded methods:
    \texttt{search(String title)} and \texttt{search(String author, int year)}.
    (Use arrays to store books; collections are not required.)

  \item \textbf{Prevent invalid states:} Identify 5 invariants from everyday entities (bank account, booking, student, inventory item, rectangle).
    For any 2 invariants, show how you would enforce them in constructors/methods (write code snippets).

  \item \textbf{Integrated demo app (packages + polymorphism):} Create a small multi-package project with at least:
    one interface, two implementing classes, one base class with one subclass, and one \texttt{main} that demonstrates dynamic dispatch.
    Provide the package tree and a short explanation of how to compile/run.
\end{enumerate}

\vfill
\noindent\textit{End of Assignment}

\end{document}

